\begin{abstract}
\textbf{Background:} Severe thalassemia can require lifelong transfusion and
iron chelation, while curative or near-curative options remain limited by
eligibility, infrastructure, and cost.
\textbf{Objective:} Nakafa Lab is building an AI-assisted, source-backed
research program to identify credible and more affordable disease-modifying
paths for thalassemia, with Indonesia access as an explicit constraint.
\textbf{Methods:} We organized clinical guidelines, regulatory sources,
ClinicalTrials.gov snapshots, PubMed evidence, chemistry identity snapshots,
de-identified case-context templates, and Islamic research guardrails into a
versioned public repository. Evidence was grouped into therapy benchmarks,
mechanism lanes, candidate gates, no-lab clinician packets, and future assay
requirements.
\textbf{Current findings:} The strongest cure-oriented benchmark remains
hematopoietic stem-cell replacement or modification, including HSCT, lentiviral
gene addition, and CRISPR-enabled fetal hemoglobin reactivation. The most
practical affordable research direction is not a single unverified supplement,
but a gated program around fetal hemoglobin, globin-chain balance, erythroid
maturation, mature red-cell safety, iron overload, and specialist-reviewed
eligibility.
\textbf{Status:} This manuscript is an exploratory evidence map and research
program proposal. It is not a systematic review, clinical guideline, treatment
recommendation, or evidence that Nakafa Lab has found a cure.
\end{abstract}
